\chapter{Reflux in Babies}

\section*{Definition}
Reflux in babies is the passage of stomach contents back into the esophagus, often causing effortless milk spit-up. It is common because the infant lower esophageal sphincter is immature and usually improves with growth.

\section*{When to See a Doctor}
Urgent evaluation is needed for green or bloody vomit, projectile vomiting, breathing difficulty, poor feeding, dehydration, abdominal distension, fever, failure to thrive, or a lethargic infant. Arrange prompt evaluation for persistent, recurrent, unexplained, or function-limiting symptoms.

\section*{How It Develops}
The symptom results from irritation, inflammation, altered organ function, injury, infection, or disruption of normal neurologic or vascular signaling. Age, pregnancy status, onset, distribution, severity, and associated findings determine the urgency and likely causes.

\section*{Causes}
\begin{itemize}
  \item Physiologic infant reflux, overfeeding, swallowing air, cow’s-milk protein allergy, gastroesophageal reflux disease, infection, obstruction, and feeding or neurologic disorders.
  \item Medication effects, environmental exposures, and underlying chronic disease may also contribute.
  \item Serious causes are less common but must be considered when symptoms are sudden, progressive, severe, or associated with systemic illness.
\end{itemize}

\section*{Related Symptoms}
\begin{itemize}
  \item Pain, fever, fatigue, nausea, weakness, or changes in normal function
  \item Bleeding, swelling, discharge, breathing difficulty, or neurologic symptoms when a serious cause is present
\end{itemize}

\section*{Medical Evaluation}
\begin{itemize}
  \item History addresses onset, duration, severity, triggers, injuries, exposures, medications, medical conditions, age, and pregnancy status when relevant.
  \item Examination includes vital signs and focused assessment of the relevant organ system, neurologic status, hydration, and function.
  \item Testing may include blood or urine studies, cultures, imaging, physiologic testing, fetal or pediatric assessment, or specialist examination.
\end{itemize}

\section*{Treatment Options}
Treatment is directed at the cause and may include supportive care, targeted medication, treatment of infection or inflammation, correction of metabolic disease, physical therapy, or specialist procedures. Persistent symptoms require reassessment.

\section*{Self-Care and Home Management}
Avoid known triggers, protect the affected area, maintain appropriate hydration, and record timing and associated symptoms. Use only age- or pregnancy-appropriate medicines recommended by a clinician. Do not delay emergency care while trying home treatment.

\section*{References}
\begin{thebibliography}{99}
\bibitem{reflux_in_babies:naspghan} Rosen R, Vandenplas Y, Singendonk M, et al. Pediatric Gastroesophageal Reflux Clinical Practice Guidelines: joint recommendations of the North American Society for Pediatric Gastroenterology, Hepatology, and Nutrition and the European Society for Pediatric Gastroenterology, Hepatology, and Nutrition. \textit{J Pediatr Gastroenterol Nutr}. 2018;66:516--554.
\bibitem{reflux_in_babies:nice} National Institute for Health and Care Excellence. Gastro-oesophageal reflux disease in children and young people: diagnosis and management. NICE guideline NG1; 2015 (updated 2019).
\bibitem{reflux_in_babies:aap} American Academy of Pediatrics. Choosing Wisely: recommendations to reduce harm — routine acid-suppressant use in infants. 2015.
\end{thebibliography}
